\documentclass[../main.tex]{subfiles}

\begin{document}

    We have been told to consider the boundary-value problem
    \[
        \frac{d^{2}y}{dx^{2}} = \cos(x),
    \]
    on the domain $[0, 2\pi]$.

    \section{Finding the true solution}\label{sec:bvp-finding-true-solution}
    First, let's find a trial solution:
    \begin{center*}
        \begin{align*}
            & \frac{d^{2}y}{dx^{2}} = \cos(x) \\
            & \implies \frac{dy}{dx} = \sin(x) + A \\
            & \implies y = -\cos(x) + Ax + B
        \end{align*}
    \end{center*}\\
    Now, using the initial conditions $y'(0) = y(2\pi) = 0$,
    \begin{center*}
        \begin{align*}
            & y'(0) = 0 = \sin(0) + A \\
            & \implies A = 0 \\
            \\
            & y(2\pi) = 0 = -\cos(2\pi) + (0 \cdot 2\pi) + B \\
            & \implies B = 1
        \end{align*}
    \end{center*}\\
    So, our true solution is $y_{*}(x) = -cos(x) + 1$.

    \section{Forming a matrix inversion problem}\label{sec:bvp-forming-a-matrix-inversion-problem}
    We start by discretising our domain $[0, 2\pi]$ into $N$ points, giving $x_{0}, x_{1}, \dots, x_{N - 1}$ and
    $\Delta x = \frac{2\pi}{N - 1}$.
    Then, using a finite difference approximation with interior points $i = 1, 2, \dots, N - 2$,
    \begin{center*}
        \begin{align*}
            & \frac{d^{2}y}{dx^{2}} \approx \frac{y_{i - 1} - 2y_{i} + y_{i + 1}}{\Delta x^{2}} \\
            & \implies \cos(x_{i}) = \frac{y_{i - 1} - 2y_{i} + y_{i + 1}}{\Delta x^{2}} \\
            & \implies \Delta x^{2} \cos(x_{i}) = y_{i - 1} - 2y_{i} + y_{i + 1}
        \end{align*}
    \end{center*}\\
    giving us the linear system
    \[
        T \textbf{y} = \Delta x^{2} \textbf{f}
    \]
    where $T$ is an $N \times N$ tridiagonal matrix, and $\textbf{y}$ and $\textbf{f}$ are vectors:
    \begin{center*}
        \begin{align*}
            \begin{pmatrix}
                1 & 0 & 0 & \dots & 0 \\
                1 & -2 & 1 & \dots & 0 \\
                0 & \ddots & \ddots & \ddots & 0 \\
                0 & \dots & 1 & -2 & 1 \\
                0 & \dots & 0 & 0 & 1
            \end{pmatrix}
            \begin{pmatrix}
                y_{0} \\
                \vdots \\
                y_{i} \\
                \vdots \\
                y_{N - 1}
            \end{pmatrix} =
            \Delta x^{2} \begin{pmatrix}
                0 \\
                \cos(x_{1}) \\
                \vdots \\
                \cos(x_{N - 2}) \\
                0
            \end{pmatrix}
        \end{align*}
    \end{center*}\\

    \subsection*{Boundary conditions}\label{subsec:bvp-boundary-conditions}
    We have two boundary conditions to consider on our domain $[0, 2\pi]$:
    \begin{itemize}
        \item Neumann boundary condition on the left boundary: \\ \\
        The Neumann condition means that we need to prescribe the derivative $y'(0) = 0$ at the boundary.
        To implement this, we approximate the derivative using a forward Euler approximation:
        \begin{center*}
            \begin{align*}
                & y'(0) \approx \frac{y_{1} - y_{0}}{\Delta x} \\
                & \implies 0 = \frac{y_{1} - y_{0}}{\Delta x} \\
                & \implies 0 = y_{1} - y_{0} \\
                & \implies y_{0} - y_{1} = 0
            \end{align*}
        \end{center*}\\
        Therefore, the first row of our matrix $T$ is $(1, -1, 0, \dots, 0)$.
        \item Dirichlet boundary condition on the right boundary: \\ \\
        The Dirichlet condition means that we need to prescribe the value of the function $y(2\pi) = 0$ at the boundary.
        Since we're working with a discretised system with the domain split into $N$ points where $x_{N - 1} = 2\pi$,
        the Dirichlet condition applies to the last point.
        \[
            \mathbf{y}_{N - 1} = 0
        \]
        Then, we replace the last row in our matrix $T$ with $(0, 0, \dots, 0, 1)$ and we ensure $\mathbf{y}_{N - 1} = 0$
        by setting $\mathbf{f}_{N - 1} = 0$.
    \end{itemize}
    These conditions give us the linear system
    \begin{center*}
        \begin{align*}
            \begin{pmatrix}
                1 & -1 & 0 & \dots & 0 \\
                1 & -2 & 1 & \dots & 0 \\
                0 & \ddots & \ddots & \ddots & 0 \\
                0 & \dots & 1 & -2 & 1 \\
                0 & \dots & 0 & 0 & 1
            \end{pmatrix}
            \begin{pmatrix}
                y_{0} \\
                \vdots \\
                y_{i} \\
                \vdots \\
                y_{N - 1}
            \end{pmatrix} =
            \Delta x^{2} \begin{pmatrix}
                0 \\
                \cos(x_{1}) \\
                \vdots \\
                \cos(x_{N - 2}) \\
                0
            \end{pmatrix}
        \end{align*}
    \end{center*}\\

    \section{Solving our linear system}\label{sec:bvp-solving-linear-system}
    To solve our linear system for $\mathbf{y}$ we rearrange $T\mathbf{y} = \Delta x^{2}\mathbf{f}$ to equal
    $\mathbf{y} = T^{-1} \Delta x^{2} \mathbf{f}$.
    To do this I used the \texttt{numpy.linalg} Python module to solve our equation.
    This gave us $N$ plotted points that estimate the solution of the boundary-value problem
    \[
        \frac{d^{2}y}{dx^{2}} = \cos(x).
    \]
    Which when placed on a graph alongside the analytical true solution $y_{*} = -\cos(x) + 1$ looks like
    \begin{figure}[H]
        \centering
        \includegraphics[width=230px]{boundary-value-problem-N-51} \\
        \textbf{Figure:} Numerical solution with $N = 51$.
        \label{fig:bvp-n-51}
    \end{figure}

    \section{The relationship between error and $\Delta x$}\label{sec:bvp-relationship-error-and-grid-length}
    The error of a numerical solution with grid length $\Delta x$ is given by
    \[
        \epsilon = \overline{|y_{num} - y_{*}|} \text{,}
    \]
    where $y_{num}$ is the numerical solution and $y_{*}$ is the true solution $y_{*} = -\cos(x) + 1$.
    Plotting the error $\epsilon$ against $\Delta x$ on a log-log scale gives:
    \begin{figure}[H]
        \centering
        \includegraphics[width=230px]{boundary-value-problem-loglog} \\
        \textbf{Figure:} Error values plotted against $\Delta x$ on a log-log scale.
        \label{fig:bvp-error-vs-grid-length}
    \end{figure}
    As you can see, the line is roughly $y = x$, and so we can say that the error in the numerical solutions
    satisfies the relation
    \[
        \epsilon \sim \Delta x^{n} \text{,}
    \]
    where $n$ is the order of convergence.
    \\ \\
    From the log-log plot, the slope is $n \approx 1$, inferring first-order convergence.
    This is what we expected to see, because:
    \begin{itemize}
        \item The Neumann boundary condition was approximated by a forward Euler scheme - a first-order scheme,
        \item The interior points use a second-order central difference for $y''$, but the error $\epsilon$
        is dominated by the lowest-order term.
    \end{itemize}
    Thus, the method exhibits \textbf{first-order accuracy}.

    \subsection*{Code used to solve the system}\label{subsec:bvp-code}
    I ran the following in my terminal (usage covered in
    \hyperref[subsubsec:bvp-appendix]{Boundary-value matrix problem} under
    \hyperref[subsec:usage-of-modules]{Usage of modules}):
    \begin{lstlisting}[style=console,label={lst:bvp}]
C:\User\CaseStudy> python main.py boundary
Order: 1.00
    \end{lstlisting}

\end{document}
